%% This is the ctufit-thesis example file. It is used to produce theses
%% for submission to Czech Technical University, Faculty of Information Technology.
%%
%% Get the newest version from
%% https://gitlab.fit.cvut.cz/theses-templates/FITthesis-LaTeX
%%
%%
%% Copyright 2021, Eliska Sestakova and Ondrej Guth
%%
%% This work may be distributed and/or modified under the
%% conditions of the LaTeX Project Public Licenese, either version 1.3
%% of this license or (at your option) any later version.
%% The latest version of this license is in
%%  https://www.latex-project.org/lppl.txt
%% and version 1.3 or later is part of all distributions of LaTeX
%% version 2005/12/01 or later.
%%
%% This work has the LPPL maintenance status `maintained'.
%%
%% The current maintainer of this work is Ondrej Guth.
%% Contact ondrej.guth@fit.cvut.cz for bug reports.
%% Alternatively, submit bug reports into the tracker at
%% https://gitlab.fit.cvut.cz/theses-templates/FITthesis-LaTeX/issues
%%
%%

%%%%%%%%%%%%%%%%%%%%%%%%%%%%%%%%%%%%%%%%%
% CLASS OPTIONS
% language: czech/english/slovak
% thesis type: bachelor/master/dissertation
% colour: bw for black&white OR no option for default colour scheme
%%%%%%%%%%%%%%%%%%%%%%%%%%%%%%%%%%%%%%%%%
\documentclass[english,bachelor,unicode]{ctufit-thesis}


%%%%%%%%%%%%%%%%%%%%%%%%%%%%%%%%%%
% FILL IN THIS INFORMATION
%%%%%%%%%%%%%%%%%%%%%%%%%%%%%%%%%%
\ctufittitle{Hiding Leaders in Covert Networks: A Computational Complexity Perspective} % replace with the title of your thesis
\ctufitauthorfull{Patrik Drbal} % replace with your full name (first name(s) and then family name(s) / surname(s)) including academic degrees
\ctufitauthorsurnames{Drbal} % replace with your surname(s) / family name(s)
\ctufitauthorgivennames{Patrik} % replace with your first name(s) / given name(s)
\ctufitsupervisor{Ing.\,Šimon Schierreich} % replace with name of your supervisor/advisor (include academic degrees)
\ctufitdepartment{Department of Theoretical Computer Science} % replace with the department of your defence
\ctufityear{2023} % replace with the year of your defence
\ctufitdeclarationplace{Prague} % replace with the place where you sign the declaration
\ctufitdeclarationdate{\today} % replace with the date of signature of the declaration
\ctufitabstractENG{
    In this thesis, we provide an introduction to the topic of covert networks and their analysis,
    with emphasis on the \HL problem
    which we study with respect to the degree centrality measure
    and using the framework of parameterized computational complexity.

    We analyze results for the \HL problem provided in the literature and put them into a perspective of the parameterized complexity.
    We also point out that there are multiple definitions of the problem
    and review some of the results from the literature with respect to the definition we have picked.

    We then bring new results to the domain by showing that
    the problem is \Wh when parameterized by $b+d$ even when $|L|=1$, and \pNPh with respect to $|L|$,
    where $b$ denotes a number of edges allowed to be added into given network,
    $d$ denotes a number of followers needed to have centrality at least as high as any of leaders
    and $L$ denotes a set of leaders.
    We also show that the problem is in \XP when parameterized by $b$ or $d$.

    In the end, we put together our results and their corollaries -- for parameters $|L|$, $b$, $d$ --
    and results from the literature -- for parameter $\lambda$, where $\lambda$ denotes a maximum degree among leaders --
    and visualize them altogether in an overviewing graph,
    providing an easily accessible summary of the current state of the research of \HL
    given the decision version of the problem, the degree centrality and the four aforementioned parameters.
}
\ctufitabstractCZE{
    V této práci poskytujeme úvod do problematiky skrytých sítí (\emph{covert networks}) a jejich analýzy,
    s důrazem na problém \HL, který zkoumáme s ohledem na míru centrality měřenou stupněm uzlu (\emph{degree centrality})
    a pomocí frameworku parametrizované výpočetní složitosti.

    Analyzujeme výsledky pro problém \HL uvedené v literatuře a dáváme je do perspektivy parametrizované složitosti.
    Zdůrazňujeme také, že existuje více definic tohoto problému a přezkoumáváme některé výsledky z literatury
    s ohledem na námi vybranou definici.

    Dále přinášíme nové výsledky a ukazujeme, že problém je \textsf{W}[1]-těžký pro parametr $b+d$ i když $|L|=1$,
    a \textsf{para-NP}-těžký vzhledem k $|L|$,
    kde $b$ označuje počet hran, které mohou být přidány do dané sítě,
    $d$ označuje požadovaný počet následovníků, kteří musí mít centralitu alespoň tak vysokou jako kterýkoli z vůdců,
    a $L$ označuje množinu vůdců.
    Také ukazujeme, že problém je v \XP při parametrizaci parametry $b$ nebo $d$.

    V závěru práce spojujeme naše výsledky a jejich důsledky -- pro parametry $|L|$, $b$, $d$ --
    s výsledky z literatury -- pro parametr $\lambda$, kde $\lambda$ označuje maximální stupeň mezi vůdci --
    a vizualizujeme je společně v přehledném grafu,
    který poskytuje snadno dostupné shrnutí současného stavu výzkumu \HL problému
    vzhledem k rozhodovací verzi tohoto problému, centralitě měřené stupněm uzlu a čtyřem výše uvedeným parametrům.
}
\ctufitkeywordsENG{Hiding Leaders problem, covert networks, parameterized complexity, degree centrality}
\ctufitkeywordsCZE{problém Hiding Leaders, covert networks, parametrizovaná složitost, degree centrality}
%%%%%%%%%%%%%%%%%%%%%%%%%%%%%%%%%%
% END FILL IN
%%%%%%%%%%%%%%%%%%%%%%%%%%%%%%%%%%


%%%%%%%%%%%%%%%%%%%%%%%%%%%%%%%%%%
% CUSTOMIZATION of this template
% Skip this part or alter it if you know what you are doing.
%%%%%%%%%%%%%%%%%%%%%%%%%%%%%%%%%%
\RequirePackage{iftex}[2020/03/06]
\iftutex % XeLaTeX and LuaLaTeX
    \RequirePackage{ellipsis}[2020/05/22] %ellipsis workaround for XeLaTeX
\else
    \RequirePackage[utf8]{inputenc}[2018/08/11] %this file encoding
    \RequirePackage{lmodern}[2009/10/30] % vector flavor of Computer Modern font
\fi

% hyperlinks
\RequirePackage[pdfpagelayout=TwoPageRight,colorlinks=false,allcolors=decoration,pdfborder={0 0 0.1}]{hyperref}[2020-05-15]

% uncomment the following to hide all hyperlinks
% \RequirePackage[pdfpagelayout=TwoPageRight,hidelinks]{hyperref}[2020-05-15]

\RequirePackage{pdfpages}[2020/01/28]

\setcounter{secnumdepth}{4} % numbering sections; 4: subsubsection
%%%%%%%%%%%%%%%%%%%%%%%%%%%%%%%%%%
% CUSTOMIZATION of this template END
%%%%%%%%%%%%%%%%%%%%%%%%%%%%%%%%%%


%%%%%%%%%%%%%%%%%%%%%%
% DEMO CONTENTS SETTINGS
% You may choose to modify this part.
%%%%%%%%%%%%%%%%%%%%%%
\usepackage{dirtree}
\usepackage{lipsum,tikz}
    \usetikzlibrary{shapes.geometric}
    \usetikzlibrary{calc}
    \usetikzlibrary{arrows.meta,bending,positioning}
\usepackage{csquotes}
\usepackage[style=iso-numeric,backend=bibtex]{biblatex}
\addbibresource{text/bib-database.bib}
\usepackage{listings} % typesetting of sources
% \usepackage{minted} % typesetting of sources
\usepackage{todonotes}
\usepackage{xspace}
\usepackage{subcaption}
\usepackage{xevlna}
\usepackage[most]{tcolorbox}

%theorems, definitions, etc.
\theoremstyle{plain}
\newtheorem{theorem}{Theorem}
\newtheorem{lemma}[theorem]{Lemma}
\newtheorem{corollary}[theorem]{Corollary}
\newtheorem{proposition}[theorem]{Proposition}
\newtheorem{definition}[theorem]{Definition}
\theoremstyle{definition}
\newtheorem{example}[theorem]{Example}
\theoremstyle{remark}
\newtheorem{note}[theorem]{Note}
\newtheorem*{note*}{Note}
\newtheorem{remark}[theorem]{Remark}
\newtheorem*{remark*}{Remark}
\numberwithin{theorem}{chapter}
%theorems, definitions, etc. END
%%%%%%%%%%%%%%%%%%%%%%
% DEMO CONTENTS SETTINGS END
%%%%%%%%%%%%%%%%%%%%%%


\begin{document} 
\frontmatter\frontmatterinit % do not remove these two commands

\includepdf[pages={1-}]{assignment-include.pdf} % replace that file with your thesis assignment provided by study office

\thispagestyle{empty}\cleardoublepage\maketitle % do not remove these three commands

\imprintpage % do not remove this command

\tableofcontents % do not remove this command


%%%%%%%%%%%%%%%%%%%%%%
% list of other contents: figures, tables, code listings, algorithms, etc.
% add/remove commands accordingly
%%%%%%%%%%%%%%%%%%%%%%
\listoffigures % list of figures
\begingroup
\let\clearpage\relax
% \listoftables % list of tables
% \lstlistoflistings % list of source code listings generated by the listings package
% \listoflistings % list of source code listings generated by the minted package
\endgroup
%%%%%%%%%%%%%%%%%%%%%%
% list of other contents END
%%%%%%%%%%%%%%%%%%%%%%


%%%%%%%%%%%%%%%%%%%
% ACKNOWLEDGMENT
% FILL IN / MODIFY
% This is a place to thank people for helping you. It is common to thank your supervisor.
%%%%%%%%%%%%%%%%%%%
\begin{acknowledgmentpage}
    On this place, I would like to express my sincere gratitude to the supervisor of this work, Ing. Šimon Schierreich,
    for generously providing me with his knowledge, expertise and time.
    Especially for giving me consultations whenever I needed, as
    they always strengthen my motivatation to keep working and I always left them a little smarter,
    for willingly and helpfully pointing out mistakes I made, and
    for being my guide to the whole process, letting me dive deeper into the fascinating world of
    theoretical computer science.

    I further wish to thank Ing. Dana Vynikarová, Ph.D. for teaching me all the formal aspects of writing
    and structuring scientific text, and for her an ever-present smile on her face
    and willingness to help.

    I am also grateful to my parents for always believing in me and supporting me.

    Last but not least, my deepest thanks belong to my girlfriend Dorotka for all the love and support she gives me.
    Without her, I would hardly be able to achieve such a peace of mind and determination that
    allowed me to properly work on this thesis.
\end{acknowledgmentpage} 
%%%%%%%%%%%%%%%%%%%
% ACKNOWLEDGMENT END
%%%%%%%%%%%%%%%%%%%


%%%%%%%%%%%%%%%%%%%
% DECLARATION
% FILL IN / MODIFY
%%%%%%%%%%%%%%%%%%%
% INSTRUCTIONS
% ENG: choose one of approved texts of the declaration. DO NOT CREATE YOUR OWN. Find the approved texts at https://courses.fit.cvut.cz/SFE/download/index.html#_documents (document Declaration for FT in English)
% CZE/SLO: Vyberte jedno z fakultou schvalenych prohlaseni. NEVKLADEJTE VLASTNI TEXT. Schvalena prohlaseni najdete zde: https://courses.fit.cvut.cz/SZZ/dokumenty/index.html#_dokumenty (prohlášení do ZP)
\begin{declarationpage}
    I hereby declare that the presented thesis is my own work and that I have cited all sources of information in accordance with the Guideline
    for adhering to ethical principles when elaborating an academic final thesis.
    I acknowledge that my thesis is subject to the rights and obligations stipulated by the Act No. 121/2000 Coll.,
    the Copyright Act, as amended, in particular that the Czech Technical University in Prague has the right to conclude
    a license agreement on the utilization of this thesis as a school work under the provisions of Article 60 (1) of the Act.
\end{declarationpage}
%%%%%%%%%%%%%%%%%%%
% DECLARATION END
%%%%%%%%%%%%%%%%%%%


\printabstractpage % do not remove this command


%%%%%%%%%%%%%%%%%%%
% ABBREVIATIONS
% FILL IN / MODIFY
% OR REMOVE ENTIRELY
% List the abbreviations in lexicography order.
%%%%%%%%%%%%%%%%%%%
\chapter{List of abbreviations}

\begin{tabular}{rl}
    CNA  & Covert network analysis\\
    FPT  & Fixed-parameter tractable\\
    HL   & Hiding Leaders\\
    SNA  & Social network analysis
\end{tabular}
%%%%%%%%%%%%%%%%%%%
% ABBREVIATIONS END
%%%%%%%%%%%%%%%%%%%


\mainmatter\mainmatterinit % do not remove these two commands


%%%%%%%%%%%%%%%%%%%
% THE THESIS
% MODIFY ANYTHING BELOW THIS LINE
%%%%%%%%%%%%%%%%%%%
%---------------------------------------------------------------
\chapter*{Introduction}\addcontentsline{toc}{chapter}{Introduction}\markboth{Introduction}{Introduction}
%---------------------------------------------------------------

\setcounter{page}{1}

Covert networks, or covert organizations, are social structures whose one of the main concerns is to operate in secret,
concealed from the view of public, the view of governments and intelligence agencies or the view of any other unwanted subject.
Covert network analysis (CNA), as an important part of social network analysis (SNA),
is then a set of tools and techniques used to study covert networks and their members,
aiming to uncover connections between the members and identifying those with a great
influence and thus importance, despite many connections remaining unknown during this identification.
Such influential users are called leaders.
One of the most used techniques in the SNA toolset, and CNA in particular, for detecting leaders of covert networks
are \emph{centrality measures},
the concept introduced by Bavelas \cite{Bavelas1948} describing how important, or central,
a member of a given network inside this network is.
Being crucial for their working, leaders try to stay undetected by such measures,
for which they use various counter-techniques such as
network modification by adding or deleting edges, or creating whole new networks from scratch.

In this thesis, we survey covert networks and hiding inside of them from a perspective of the \HL problem (\HLshort), first defined by
Waniek et al. \cite{Waniek2017} and further studied by various authors \cite{Dey2019,Waniek2021,Mohan2023}.
More concretely, our main concern is the \HL problem with respect to the degree centrality measure
and using the framework of parameterized computational complexity.
Because the problem is \NPh by itself, we expose it to various parameters and
see how these parameters affect its computational complexity.
Our another goal is to analyze the problem variants studied in the literature
and describe the current state of the research in the language of the parameterized complexity framework.
We provide a profound description of the \HL problem in Chapter \ref{chapter:ProblemStatement}, for introduction to
the theory of parameterized complexity, please refer to Section \ref{section:ParamComp}.

%---------------------------------------------------------------
\chapter{Literature review}
%---------------------------------------------------------------

In this chapter, we bring a thorough examination of the literature on two topics;
social network analysis together with covert networks, and parameterized complexity.
The study of the former topic provides useful context for the \HL problem, whereas
the latter topic gives us tools to better reason about \NPh problems, among which, as the reader will see, the \HL problem surely belongs.

%---------------------------------------------------------------
\section{Social network analysis \& covert networks}

First of all, we mention the publication from Morselli \cite{Morselli2009} that the reader can refer to get
an extensive study of criminal networks.
Next, for a recent and exhaustive overview of SNA techniques used for covert network disruption,
the reader can use the paper from Ficara et al. \cite{Ficara2022}.

Arguably the most extensive line of research is that in which authors use covert networks to analyze structures consisting of
harmful actors.
That can be terrorist organizations or other structures interested in involving in illegal activities. These networks are
also called \emph{dark networks}~\cite{Raab2003}.
Indeed, research in this direction has a positive impact on understanding such harmful organizations,
shedding light on their inner processes and the most potentially hazardous actors and helping in dismantling the entire network,
preventing further operations and minimizing the damage done.
The understanding of how a covert network works inside can lead to an understanding of how it can be broken.
Among many others, we mention the work of these authors \cite{Waniek2017,Dey2019,Raab2003,Lindelauf2009,Xu2005,Ressler2006, SaavedraNieves2023}.
However, as noted by Saavedra-Nieves and Casas-Méndez \cite{SaavedraNieves2023}, this view on covert networks as organizations with harmful
intends is not the only possible since covert populations can consist for example of drug users, illegal immigrants, persecuted jews,
people with infectious diseases, activists, or ravers \cite{Oliver2014}.
The usefulness of covert network analysis then lingers in analyzing even such populations.
CNA tools, together with other methods of SNA, can be then used even, for example, for analyzing sports teams \cite{Buldú2019}.

The centrality measurement techniques mentioned in the thesis introduction are widely used because
finding the most important vertices within a given network is a natural approach when
studying real-world social networks \cite{Crescenzi2016}.
The most classical centrality measures are degree \cite{Shaw1954}, closeness \cite{Beauchamp1965},
betweenness \cite{Anthonisse1971,Freeman1977}, and core \cite{Seidman1983} centrality.
On top of the standard centrality measures, there are more recent and advanced techniques suitable for measuring one's rank within a network,
for example, the \emph{Game of Thieves} algorithm from Mocanu et al. \cite{Mocanu2018}, which can be used even within massive networks.
For the sake of completeness, we also mention that there also exist centrality measures on edges, e.g.,
\emph{WERW-Kpath} algorithm from {De Meo} \cite{DeMeo2013}.
The two aforementioned algorithms are more reviewed by Ficara et al. \cite{Ficara2021}.

After detecting influential members within a given network using centrality measures,
it is often desirable to talk about their influence in these networks, for which
various \emph{models of influence} can be used.
Frequently used models include \emph{Independent Cascade} model and \emph{Linear Threshold} model.
More on this topic can be found in the book from Easley and Kleinberg \cite{Easley2010}.

Apart from the centrality measures, there are other important tools SNA has to offer.
One of them are \emph{node similarity measures} which are with advantage used in \emph{link prediction} \cite{Zhou2009,Wang2015}.
Link prediction also plays an important part in the research of covert networks
as it allows to predict the existence of otherwise hidden connections between members, potentially leading to
identification of important members and relations between them.
This is why covert network users might want to manipulate the network to minimize the efficiency of such techniques \cite{Zhou2019}.
Besides CNA, link prediction has other useful applications for example in recommendation systems~\cite{Huang2005,Talasu2017}.

Next, we mention that leaders of covert networks often (but not always \cite{Fatih2012}) face what is called
the \emph{efficiency/security trade-off} \cite{Morselli2007} --
the dilemma between staying sufficiently hidden, while keeping enough influence over the network --
which is the often-mentioned topic in the literature about covert networks \cite{Waniek2017,Lindelauf2009,Crossley2012}.
However, the ways in which this trade-off is addressed vary in the literature.
Waniek et al. \cite{Waniek2017}, for example, approach this by modeling the security from the perspective of
centrality measures (centrality measure based secrecy), whereas the efficiency from the perspective of models of influence.
They then show how to construct a network designed specifically to hide its leaders, while
keeping their ability to influence the rest of the network.
Lindelauf et al. \cite{Lindelauf2009}, on the other hand, combine graph and game theory to model this dilemma.

When studying covert networks, we often talk about two parties, the \emph{evaders} and the \emph{seekers};
evaders are members (often leaders) of some covert network;
seekers are then those who try to identify the evaders.
Evaders are typically aware of the importance of their network structure, which can be seen from
the efficiency/security dilemma.
However, when talking about seekers and evaders, authors typically consider these three scenarios:
(i) Evaders act unstrategically in the sense that they are unaware of
the existence of seekers and SNA tools seekers use to reveal them;
(ii) Evaders behave as strategic actors, well-awared of the centrality analysis
seekers use against them \cite{Waniek2017,Dey2019,Waniek2016,Dey2020}.
(iii) Both seekers and evaders are strategic.
To this end, we mention the work from Waniek et al. \cite{Waniek2021} in which the authors propose a strategy
which tells the seekers which centrality measures they should use
to maximize the chances of detecting a leader of a covert network.

In the terminology of this section, the \HL problem describes a covert network consisting of strategic leaders
in the role of evaders, which are trying to evade being detected by the use of centrality measuring techniques,
for which they modify their network by adding new connections between network members.
The leaders do not delete any edges because they face the efficiency/security trade-off and
deleting edges may decrease their efficiency in the network.
The problem is not bound with any influence model and hence can be studied independently of it,
leaving the ``efficiency'' part of the efficiency/security trade-off untouched.


%---------------------------------------------------------------
\section{Parameterized computational complexity}\label{section:litRev:ParamCompl}

What follows is a summary of the literature on the topic of parameterized complexity,
together with a brief introduction to the topic.
For a thorough introduction to the topic, please refer to Section \ref{section:ParamComp}.

The classical complexity analysis is often insufficient when dealing with \NPh problems,
which are problems that are computationally intensive to solve --
they are so-called \emph{intractable}.
The main reason the classical complexity analysis is an insufficient tool to deal with \NPh problems is that
it cannot distinguish \NPh problems that are in some sense ``harder'' than others,
or distinguish between ``hard'' and ``easy'' instances of some \NPh problem.
This inability of the classical approach to distinguish between the hardness of various \NPh problems and their instances
led to the development of the parameterized complexity theory.
Parameterized complexity is a direct generalization of the classical complexity theory and it arms its users
with the ability to analyze the running time of algorithms, and thus computational complexities of problems
they try to solve, in finer detail.

The foundations of parameterized complexity were laid by Downey and Fellows
in the series of papers from years 1992 to 1995 \cite{Downey1995.1,Downey1995.2,Downey1993,Downey1995.4},
which the authors further presented later in 1999 in their book \cite{Downey1999},
which was refined and once more published in 2013 \cite{Downey2013}.
Other relevant and potentially useful literature on this topic are
two books by Flum with Grohe \cite{Flum2006} and by Niedermeier \cite{Niedermeier2006} from 2006, book by Hans et al. \cite{Hans2012} from 2012,
book by Cygan et al. \cite{Cygan2015} from 2015 (which to a great extent covers knowledge from the previous literature)
and two books by Haan \cite{Haan2019} and by Fomin et al. (focused mainly on kernelization) \cite{Fomin2019} from 2019.
However, keep in mind that there is a highly active line of research in the field of parameterized complexity, so,
as researchers keep making new discoveries, information presented in the work of mentioned authors may not necessarily be up-to-date.

%---------------------------------------------------------------
\chapter{Preliminaries}
%---------------------------------------------------------------

Before proceeding to the main parts of the work,
let us begin with a couple of definitions and theory the reader might find useful to fully understand later chapters.
This chapter introduces preliminaries on graph theory, including graph problems and centrality measures, and 
classical computational complexity followed by parameterized complexity.


%---------------------------------------------------------------
\section{Graph theory}

This section introduces basic concepts from graph theory.
For a thorough revision on this topic, please refer to the monograph of Diestel \cite{Diestel2018}.

% Graph
\begin{definition}[Graph]
    \emph{Graph} $G$ is an ordered pair of disjoint sets $V$ and $E$, $G=(V,E)$.
    Set $V$ is a finite set of arbitrary items called vertices (or nodes).
    Set $E$ is a finite set of unordered pairs of vertices called edges, $E \subseteq \{\{x,y\} \mid x,y \in V \wedge x \neq y\}$.
    Sets of all vertices $V$ and edges $E$ of graph $G$ are in general denoted by $V(G)$ and $E(G)$ respectively.
\end{definition}

We use a standard graph definition, where graphs are undirected (edges are not ordered),
unweighted (edges has uniformed ``value''), without loops (edges are pairs of distinct vertices) and
without multiple edges (edges are unique pairs).
These graphs are called \emph{simple} \cite[p.~29]{Diestel2018}.

In this thesis, the terms ``graph'' and ``network'' are interchangeable as we only use graphs to describe networks.
By network we mean such graph, where vertices represents members of the network 
and edges represents, in some sense, a connection between two given members, e.g., enabling their communication.

% Degree
\begin{definition}[Vertex degree]
    A degree of vertex $v$ in some graph $G$, denoted as $\deg(v)$, is the number of neighbors of $v$.
    In other words, it is the number of edges vertex $v$ is a part of, i.e.,
    $\Big|\{x \in V(G) \vert \{v,x\} \in E(G) \vee \{x,v\} \in E(G)\}\Big|$.
\end{definition}

A set of neighbors of some vertex $v$ is often denoted as $N(v)$, so then $\deg(v) = |N(v)|$.

% r-regular graph
\begin{definition}[$r$-regular graph]
    A graph $G$ is \emph{$r$-regular} if the degree of each vertex from $V(G)$ is $r$, i.e., $\forall v \in V(G) \colon \deg(v) = r$.
\end{definition}

% Regular graph
\begin{definition}[Regular graph]
    A graph $G$ is \emph{regular} if there exists some $r \in \mathbb{N}$ for which~$G$ is $r$-regular.
\end{definition}

% Complement graph
\begin{definition}[Complement graph]
    A \emph{complement graph} $\overline{G}$ of graph $G$ is the graph on vertices $V(G)$,
    where two vertices are adjacent (there is an edge between them) if and only if
    they are not adjacent in $G$, i.e., $\forall e \in V(G) \times V(G) \colon e \in E(\overline{G}) \Leftrightarrow e \notin E(G)$.
\end{definition}

% Induced subgraph
\begin{definition}[Induced subgraph]
    An \emph{induced subgraph} $G[S]$ of graph $G$ is the graph on vertices $V(G[S]) = S \subset V(G)$,
    where two vertices are adjacent if and only if they are adjacent in $G$, i.e.,
    $(\forall u,v \in S)(\{u,v\} \in E(G[S]) \Leftrightarrow \{u,v\} \in E(G))$.
\end{definition}

% Complete graph
\begin{definition}[Complete graph]
    A \emph{complete graph} is the graph $G$ if there is an edge between every pair of vertices from $V(G)$, i.e.,
    $(\forall u,v \in V(G))(\{u,v\} \in E(G))$.
\end{definition}

% Clique
\begin{definition}[Clique]
    A clique $C$ is the subgraph of graph $G$ where $G[E(C)]$ is a complete graph.
\end{definition}

% Vertex cover
\begin{definition}[Vertex cover]
    A \emph{vertex cover} of graph $G$ is the subset $S \subseteq V(G)$ where each edge from $E(G)$ is
    covered by some vertex from $S$. In other words, at least one endpoint of each edge is present in $S$, i.e.,
    $(\forall \{u,v\} \in E(G))(u \in S \vee v \in S)$.
\end{definition}


\subsection{Graph problems}

What follows is a definition of the \NPh graph problem we will use in Proof \ref{proofDB}.

% k-Clique
\begin{definition}[$k$-\textsc{Clique} problem]
    Given graph $G$, the $k$-\textsc{Clique} problem is to determine if
    there exists a clique $C$ in $G$ where $|V(C)| = k$.
\end{definition}

There are many other graph problems related to cliques. The $k$-\textsc{Clique} problem is a decision problem but
there are also optimization variants, for example the \textsc{Maximum Clique} problem.


\subsection{Centrality measures}

Centrality measure, or centrality for short, is a measure from graph theory which is widely used in social network analysis.
Generally, centrality is a function $c: G \times V \rightarrow \mathbb{R}$ which describes the importance of a node in given network.
With centrality, we can measure a ranking (position among other nodes) of given node within their network.
What follows is a definition of a degree centrality introduced by Shaw \cite{Shaw1954}, but there are other centralities used in SNA like
closeness~\cite{Beauchamp1965}, betweenness \cite{Anthonisse1971,Freeman1977} or core \cite{Seidman1983} centrality.

\begin{definition}[Degree centrality]
    A degree centrality measures an importance of a vertex by its degree.
    A degree centrality of the vertex $v$ in network $G$ is defined as:
    $$c_{deg}(G, v) = \deg(v).$$
\end{definition}


%---------------------------------------------------------------
\section{Classical computational complexity}

In this section, we briefly describe the two most fundamental classical complexity classes, \Po and \NP.
Detailed revision on this topic can be found in the textbook from Arora and Barak \cite{Arora2009}.

Let us start with defining the \textsf{DTIME} complexity class, which we then use to define the \Po class.

\begin{definition}[\textsf{DTIME}]
    A complexity class \textsf{DTIME}(f(n)) is the set of all decision problems that are computable in time
    $c \cdot f(n)$ for some constant $c > 0$ and some function $f \colon \mathbb{N} \rightarrow \mathbb{N}$.
\end{definition}

% P
\begin{definition}[\Po]
    A complexity class \Po is the set of all decision problems that are in class \textsf{DTIME}($n^c$)
    for any $c \geq 1$, i.e.
    $\cup_{c \geq 1} \textsf{DTIME}(n^c)$.
\end{definition}

Informally, complexity class \Po consists of all decision problems that can be solved in polynomial time.
Problems for which exists in practice efficient enough algorithm are called \emph{tractable}.

% NP
\begin{definition}[\NP]
    A complexity class \NP is the set of all decision problems $L$ for which there exist
    a polynomial $p \colon \mathbb{N} \rightarrow \mathbb{N}$ and a deterministic Turing machine $M$
    such that for every instance $x$ of $L$ $x$ is a $yes$-instance if and only if there exists
    a polynomial-size solution $u$ of $x$ that can be verified on $M$ in polynomial time.
    Such solution $u$ is then called the \emph{certificate} or the \emph{witness}.
\end{definition}

Informally, complexity class \NP consists of all decision problems that can be solved in polynomial time
when multiple steps can be done in parallel.
Note that $\Po \subseteq \NP$, because, for any problem from \Po, it takes only a polynomial amount of time to solve
it in the first place.

At the end of this section, let us briefly recall computational complexity terms \emph{hardness} and \emph{completeness}.
A problem $p$ is, given some complexity class \textsf{C}, \textsf{C-hard} if for any problem $c$ from~\textsf{C}
there exists a polynomial reduction from $c$ to $p$.
A problem $p$ is then \textsf{C-complete} if it is \textsf{C-hard} and it belongs to \textsf{C} at the same time.


%---------------------------------------------------------------
\section{Parameterized computational complexity}\label{section:ParamComp}

We use the textbook by Cygan et al. \cite{Cygan2015} as our main source of information on this topic since it
provides arguably the most complex and comprehensive overview of the theory of parameterized complexity and,
in many cases, presents the state of the art in the field.

The classical complexity analysis is often insufficient when dealing with \NPh problems.
The main reason is that there are \NPh problems that are in some sense harder than the others.
This inability of the classical approach to distinguish between the hardness of various \NPh problems
led to a development of the parameterized complexity theory.
Parameterized complexity is a direct generalization of the classical complexity theory and it arms its users
with the ability to analyze running time of algorithms, and thus computational complexities of problems
they try to solve, in a finer detail.
It gives us tools which help us to distinguish between different difficulty levels in places where
\NP-hardness fails to do so.
To achieve this, parameter complexity introduces a notion of parameterization of the input instance, where
the parameter is just some secondary measurement of the input instance.
Such parameter is then used together with the input instance size when describing a computational complexity of \NPh problems.
Given a problem, there is typically many parameters we can parameterize the problem by.
This leads us to the importance of choosing the right parameter.

When parameterizing some problem, we typically look for parameters which are small in real-world applications,
as one of the main goal when studying parameterized problems is to find algorithms for these problems in which
the running time is exponential only in the parameter, leaving the running time polynomial in the, potentially enormous, input size.
However, not every choice of the parameter let us design such an efficient algorithm, as it seems that many problems with
certain parameters admit no efficient algorithm at all.
That means the hardness, or tractability, of the problem depends on the parameter we use.
Typically, the more information a parameter carries about the input instance, the higher are chances for us to exploit this
information and design a faster algorithm.

The typical parameter we often try is the size, or other property, of the solution we are looking for.
The other type of parameter is a measure of some property of the input instance.
For example the maximum degree, regularity or treewidth measurings of the input graph.
For non-graph instances, it may be the maximum length of a string when the input instance consists of a set of strings,
or a number of variables when working with Boolean formulas.

Given problem parameterized by $p$, algorithm designers commonly talk about some function of $p$, $f(p)$.
Let us note that it is also possible to use more than one parameter.
Having parameters $k$ and $l$, we then talk about a function of $k$ and $l$, $f(k,l)$.
However, we can (and we actually do) express the parameterization by $k$ and $l$ by using just one parameter $k+l$, $f(k+l)$.

The two most fundamental complexity classes of this theory, used to reason about the complexity of parameterized problems,
are \FPT and \XP.
To distinguish between \NPh problems that are in \XP but not in \FPT, we can use another set of complexity classes, the \W-hierarchy.
There is also a complexity class called \pNP.
What follows are the formal definitions of some of these classes, together with definitions of parameterized problem
and parameterized reduction.
The hierarchy of the classes can be seen in Figure \ref{fig:complexityClasses}.


% Figure - complexity classes
\begin{figure}[t]
    \centering
    \begin{tikzpicture}
        \pgftransformscale{.8}
        %% Horizontal bar
        \draw[very thick] (6,0) -- (-6,0);
        % FPT
        \draw (-2,0) parabola bend (0,2) (2,0);
        \node at (0,1) {\FPT};
        % W[1]
        \draw (-2.5,0) parabola bend (0,3) (2.5,0);
        \node at (0,2.5) {\textsf{W}[1]};
        % W[2]
        \draw (-3,0) parabola bend (0,4) (3,0);
        \node at (0,3.5) {\textsf{W}[2]};
        % % W[t]
        \draw (-4,0) parabola bend (0,6) (4,0);
        \node[anchor=north] at (0,5.8) {
            \begin{tabular}{c}
                \textsf{W}[$t$] \\
                $\vdots$ \\
            \end{tabular}
        };
        % XP
        \draw[rounded corners] (-6,0) .. controls (-1,8) and (9,17) .. (6,0);
        \node at (4,8) {\XP};
        % para-NP
        \draw[rounded corners] (6,0) .. controls (1,8) and (-9,17) .. (-6,0);
        \node at (-4,8) {\pNP};
    \end{tikzpicture}
    \caption
    [Relations among the parameterized complexity classes]
    {
        Relations among the parameterized complexity classes.
        Inspired by the depiction from Flum and Grohe \cite[p.~97]{Flum2006}.
    }
    \label{fig:complexityClasses}
\end{figure}


% Parameterized problem
\begin{definition}[Parameterized problem]
    A parameterized problem is a language $L \subseteq \Sigma^* \times \mathbb{N}$, where
    $\Sigma$ is a fixed, finite alphabet and $\Sigma^*$ is a set of all strings over $\Sigma$.
    For an instance $(x, k) \in \Sigma^* \times \mathbb{N}$, $k$ is called the parameter.
\end{definition}

One particular problem parameterized by different parameters leads to different parameterized problems.

% Parameterized reduction
\begin{definition}[Parameterized reduction]
    Let $A,B \subseteq \Sigma^* \times \mathbb{N}$ be two parameterized problems.
    A parameterized reduction from $A$ to $B$ is an algorithm that, given an instance $(x, k)$ of $A$,
    outputs an instance $(x', k')$ of $B$ such that
    \begin{itemize}
        \item $(x, k)$ is a $yes$-instance of $A$ if and only if $(x', k')$ is a $yes$-instance of $B$,
        \item $k' \leq g(k)$ for some computable, nondecreasing function $g$, and
        \item the running time is $f(k) \cdot |x|^{\mathcal{O}(1)}$ for some computable, nondecreasing function $f$.
    \end{itemize}
\end{definition}

% FPT problem and class
\begin{definition}[FPT]
    A parameterized problem $L$ is called fixed-parameter tractable (FPT) if there exists
    \begin{itemize}
        \item an algorithm $\mathcal{A}$, called fixed-parameter algorithm, or FPT algorithm,
        \item a computable, nondecreasing function $f : \mathbb{N} \times \mathbb{N}$
        \item and a constant $c$
    \end{itemize}
    such that, given $(x,k) \in \Sigma^* \times \mathbb{N}$,
    the algorithm $\mathcal{A}$ correctly decides whether $(x, k) \in L$ in time bounded by
    $f(k) \cdot |(x,k)|^c$.
\end{definition}

The complexity class containing all fixed-parameter tractable problems is called \FPT.
An example of a problem from \FPT is $n$-variable \textsc{SAT} parameterized by $n$.
Indeed, the problem can be solved in time $\Theta(2^n)$ by simply trying each of the $2^n$
possible evaluations for the variables and checking if the evaluation is correct in time $\Theta(n)$.
Moreover, we know from the Exponential Time Hypothesis, formulated by Impagliazzo and Paturi \cite{Impagliazzo1999},
that the problem cannot be solved in time $2^{o(n)}$

The typical goal when designing FPT algorithms is to make factor $f(k)$ and constant exponent~$c$
in the running time boundary as small as possible.
We can see that if the parameter is equal to the input size, then \FPT become exponential and,
on the other hand if $k = 1$, then \FPT become~\Po~\cite{Koutensky2020}.

The only difference between FPT and para-NP problems is that algorithms for para-NP problems are
nondeterministic in general.
We do not provide the exact definition of para-NP problems as they are not important to us.
From the \pNP-hardness perspective, we can say that a problem is \pNPh if it is \NPh already for a constant value of the parameter.

% XP
\begin{definition}[XP]
    A parameterized problem $L$ is called slice-wise polynomial (XP) if there exists
    \begin{itemize}
        \item an algorithm $\mathcal{A}$
        \item and two computable, nondecreasing functions $f,g : \mathbb{N} \times \mathbb{N}$
    \end{itemize}
    such that, given $(x,k) \in \Sigma^* \times \mathbb{N}$,
    the algorithm $\mathcal{A}$ correctly decides whether $(x, k) \in L$ in time bounded by
    $f(k) \cdot |(x,k)|^{g(k)}$.
\end{definition}

The complexity class containing all slice-wise polynomial problems is called \XP.
An example of a problem from \XP is $k$-\textsc{Independent Set}.

Note that definitions of parameterized problems FTP and XP differs only in a running time boundary of algorithm $\mathcal{A}$
and so $\FPT \subseteq \XP$.
Although both FPT and XP algorithms run in polynomial time for every fixed value of the parameter,
the underlying difference between them is that FPT algorithms have the combinatorial explosion
restricticted to the parameter only, leaving the exponent of the instance size constant, in other words,
FPT algorithms are more efficient than XP algorithms.


\subsection{\W-hierarchy}

As stated by Cygan et al. \cite[p.~423]{Cygan2015},
there are thousands of natural problems which are \NPc and which are reducible to each other,
meaning that in this sense, they are equally hard and thus we can say that
they occupy the same level of hardness.
However, we cannot say the same when talking about parameterized problems as it seems there are
different levels of hardness for such problems and, in this sense, even basic problems
seem to be differently hard as they occupy different hardness levels.
For this reason, \W-hierarchy was introduced by Downey and Fellows \cite{Downey1999} as
an attempt to shed a light on these apparent differences in hardness of parameterized problems. 

The levels of \W-hierarchy are marked as \textsf{W}[$t$] for $t \in \mathbb{N} \wedge t \geq 0$,
where each level represents its own complexity class.
The most important level for us will be \Wone,
as well as the fact that the $k$-\textsc{Clique} problem is \Wone-complete \cite{Downey1999}.

We do not provide the exact definition and further description of \W-hierarchy
because it is not important for the purposes of this work.
On the other hand, what is important is a fact that we interpret the \W-hardness as an evidence that
a problem is not fixed-parameter tractable.
This interpretation is based on a general assumption that $\FPT \neq \Wone$.
To show that certain problem is \Wh, it is typically shown a parameterized reduction from some other problem,
which is already known to be \Wh, to it.

%---------------------------------------------------------------
\chapter{Problem statement}\label{chapter:ProblemStatement}
%---------------------------------------------------------------

Now we define the central problem of this thesis, the \HL problem.
We also provide a detailed description of the problem, a motivation behind its formulation,
an explanation of its hardness, an introduction to notation we use and a presentation of results from the literature. 

\begin{center}
    \line(1,0){\textwidth}
\end{center}
\begin{definition}[Hiding Leaders]\label{HL}
    Let $(G, L, b, c, d)$ be the problem instance, then
    \begin{description}
        \item $G = (V, E)$ is a network,
        \item $L \subseteq V$ are leaders and the remaining vertices $F = V \setminus L$ are followers,
        \item $c : G \times V \rightarrow \mathbb{R}$ is a centrality measure,
        \item $b$ is a maximum number of edges that we are allowed to add in $G$,
        \item $d$ is a safety margin -- a number of followers whose final centrality should be at least as high as of any leader.
    \end{description}
    Given this instance, the goal is to identify a set of maximum $b$ edges between followers W $\subseteq$ F {\texttimes} F
    such that the resulting network $G' = (V, E \cup |W|)$ contains at least $d$ followers $F' \subseteq F$
    whose centrality must be at least as high as the centrality of any leader, that is
    $$|W| \leq b$$
    and
    $$\exists_{F' \subseteq F} |F'| \geq d \wedge \forall_{f \in F'} \forall_{l \in L} c(G', f) \geq c(G', l).$$
\end{definition}
\begin{center}
    \line(1,0){\textwidth}
\end{center}

In other words, we want to \emph{``identify a set of edges to be added between the followers so that
the ranking of the leaders (based on some centrality measure) drops below a certain threshold''} \cite{Waniek2017}.

We survey \HLshort only with respect to the degree centrality measure, that is, with $c = c_{deg}$.
We denote this version of the problem as \HLdeg.
We also use symbol $\lambda$ to denote a maximum degree among leaders, this is the minimum degree that
all followers from $F'$ have to reach.
Lastly, we use symbol $\hat{A}$ to denote a set of edges between followers that are not present in $G$ and thus can be added,
$\hat{A} = E(\overline{G[F]})$.

There are two definitions of the \HL problem in the literature and they differ from each other in some details.
Namely, the definition from Waniek et al. \cite{Waniek2017} allows leaders $L$ to be equal to $V$, $L \subseteq V$,
which is a small difference since for $L = V$ the problem is trivial, but the more important difference is that
it requires centrality measure of followers $F'$ to be strictly greater than that of any leader,
$\forall_{f \in F'} \forall_{l \in L} c(G', f) > c(G', l)$.
We decided to stick to the definition from Dey \cite{Dey2019} because it seems more natural to us to let followers
to be part of a solution as soon as their centrality is at least the same as of all the leaders --
if no leader has centrality measure greater than any follower from $F'$, then
we consider the leaders hidden and safe from detection.
On the other hand, the name ``\HL'' comes from the original definition by Waniek et al. as
the name is more indicative of the possibility that there are many leaders in the network.
Later in Section \ref{section:proofRevision}, we show that complexity results for degree centrality from Waniek et al.
are correct even with this slighty different definition of \HLshort.

The motivation behind only adding edges, with no deletion involved, is that we want leaders to maintain
their existing control over the network.
Indeed, deleting edges incident to the leaders would decrease their number of connections and thus network control,
whereas deleting edges between followers would only decrease their centrality \todo{explain why},
which is the exact opposite of what we try to achieve.
Because no connections are deleted in the resulting network, all leaders keep their existing control over the network,
having increased their security but keeping their efficiency, although, new members with great influence may arise.

Also note that decreasing the value of degree centrality of any given member is a straightforward task
(in contrary with decreasing other centralities) as it only consists in cutting edges \cite{Waniek2016}.
On the other hand, \HLdeg is \NPc \cite{Waniek2017}.
The core difference between these two problems is that \HLdeg, and \HLshort in general,
is not about decreasing centrality of given member but rather about decreasing their ranking,
the relative position among other nodes with respect to the centrality measure.
To decrease member's ranking, we must increase centrality of some other members.
This fact, together with budget constraint $b$ and safety margin constraint $d$,
is what stands behind the hardness of the \HL problem, because,
as Waniek et al. \cite{Waniek2017} shown, and as we will also show in Chapter \ref{chapter:contribution},
there exist instances of \HLshort where finding a solution corresponds to solving certain \NPh problem.

From the work of Dey et al. \cite{Dey2019}, we know that \HLdeg is polynomial time solvable
if the degree of every leader is bounded by some constant.
That means that \HLdeg admits a \FPT algorithm when parameterized by degree of leaders.
Moreover, the authors present a $2$-approximation algorithm for \HLdeg which optimizes
the number of edges added.
They complement that by proving that if there exists a $(2-\epsilon)$-approximation algorithm
for the above problem for any constant $0 < \epsilon < 1$, then there exists
a $(\frac{\epsilon}{2})$-approximation algorithm for the \textsc{Densest} $k$-\textsc{Subgraph} problem.

%---------------------------------------------------------------
\chapter{Contribution}\label{chapter:contribution}
%---------------------------------------------------------------

This is the main chapter of the thesis.
On this place, we present our contribution to the topic of the \HL problem.
We first review a proof of a theorem about \HLdeg from Waniek et al. \cite{Waniek2017} and show
that the proof can be adapted for our definition of \HL.
Then, we present our own proofs for \HLdeg regarding various parameters and complexity classes.
Lastly, we put our results in the context of the work of other authors,
visualizing the parameterized complexity results for \HLdeg discovered both by us and by authors
from the literature.


%---------------------------------------------------------------
\section{Proof revision}\label{section:proofRevision}

Waniek et al. \cite{Waniek2021full} proposed three hardness results for the \HL problem given degree, closeness and betweenness
centralities. In this section, we review the first aforementioned result with respect to our definition of \HLshort.
The original theorem reads:

\begin{theorem}
    The \HL problem is \NPc given the degree centrality.
\end{theorem}

The proof that follows is almost identical to the original proof from Waniek et al. \cite{Waniek2021full},
with only small adjustments done.
The adjustments are pretty straightforward but to better understand them,
we start the proof by briefly describing the idea of the original proof.

\begin{proof}
    A common way of proving \NP-hardness is to propose a polynomial reduction from another \NPh problem to the
    problem we are proving \NP-hardness of.
    In the original proof, the authors show a reduction from the $k$-\textsc{Clique} problem defined on graph $G$.
    The idea is following: In a polynomial time, construct a new graph $H$ from $G$ in a way that
    in order to solve \HL in $H$ we must solve $k$-\textsc{Clique} in $G$ in the first place.
    To achieve that, the authors, as a starting point of $H$, first create a complement of $G$ 
    where all vertices are marked as followers. These will be the only followers in the final $H$.
    This initial construction assures that every possible solution of \HLshort will only consist of edges from $G$, that is,
    $W \subseteq E(G)$ (and $F' \subseteq V(G)$).
    Next, they add two types of leaders into $H$, each type with a different purpose.
    Leaders of the first type act like ``regularizators'' as their only job is to connect with followers (in some way) and,
    by doing so, make the degree of every follower the same and equal to $n$.
    These leaders are denoted as $X$ and $y$ in the original proof.
    The task of the second type, having no connection with followers at all,
    is then to set a maximum degree across all leaders, which followers from $F'$ will have to reach.
    These leaders are denoted as $L'$ and are practically the only concern of our adjustments to which we now proceed.

    We allow equality of centralities between followers from $F'$ and leaders in our definition of \HL.
    To compensate it, we just need to create one extra leader of the second type, that is,
    $|L'| = n+k$ as opposed of $|L'| = n+k-1$ in the original proof,
    leading to all leaders from $L'$ having degree of $n+k-1$, as opposed of $n+k-2$ in the original proof.
    Now note that this is the highest degree among the leaders and that all the followers have degree $n$,
    so, in order to be part of $F'$, they must obtain at least $(n+k-1) - n = k-1$ new neighbors.
    We got into exactly the same situation as in the original proof and since, as stated above,
    leaders from $L'$ do not have any other purpose than setting the maximum degree of all leaders,
    the rest of the arguments presented in the proof remain unchanged.

    We have just shown that allowing equality of centralities between followers from $F'$ and leaders does not
    break the proof presented by Waniek et al. \cite{Waniek2017}.
    The other change in our definition of \HL is that we allow all members to be leaders, $L = V$,
    but since such instances are trivial as they are $yes$-instances if and only if $d = 0$,
    there is nothing to change in the proof.
\end{proof}


%---------------------------------------------------------------
\section{Our results}\label{section:OurResults}

The next proof is built on the same idea Waniek et al. \cite{Waniek2017} presented in their proof of
\NP-completeness of \HLshort for the degree centrality measure.
We reviewed the proof from Waniek et al. \cite{Waniek2017} and its core ideas in Section \ref{section:proofRevision}.
We present a proof of a stronger theorem, with a simpler reduction graph and in the language of parameterized complexity.
These results are possible mainly because of a stronger assumption in the proof where we assume a regular graph as the input.

The core idea of the proof is that we construct a network in such a way that finding a solution~$W$ of the \HL problem
correspond to finding a solution to the problem we started with, essentially having the starting network
present in the constructed network in some manner and then finding a solution inside of it.

\begin{theorem}\label{theorem:DB}
    \HL parameterized by $b + d$ is \Wh even if $|L| = 1$.
\end{theorem}

\begin{proof}\label{proofDB}
    To prove this theorem, we give a parameterized reduction from the $k$-\textsc{Clique} problem on regular graphs,
    which is known to be \Wh with respect to $k$ \cite{Mathieson2008}.
    
    Suppose we have a parameterized $k$-\textsc{Clique} instance $(G, k)$, where
    $|V(G)|=n$, $G$ is a $r$-regular graph and $n-2 \geq r \geq k \geq 3$ (for $r > n-2$ or $r < k$ or $k < 3$ the problem is trivial).
    Next, let us take a parameterized problem $(\mathcal{I}, k')$, where
    $\mathcal{I} = (H, \{\ell\}, \frac{k\cdot(k-1)}{2}, c_{deg}, k)$ is an instance of \HLdeg and
    $k' = b + d = \frac{k\cdot(k-1)}{2} + k$ is the parameter.
    Graph $H$ is then constructed from graph $G$ in the following way:

    \begin{enumerate}
        \item Start with graph $H$ containing one vertex $\ell$, \emph{i.e.}, $H = (\{l\}, \{\})$.
        \item Add all vertices from $G$, \emph{i.e.}, $V(H) \leftarrow V(H) \cup V(G)$, mark these vertices as $F_1$.
        \item Add all edges from $\overline{G}$, \emph{i.e.}, $E(H) \leftarrow E(\overline{G})$.
        \item Add an edge between vertex $\ell$ and $n - r + (k - 1) = \lambda$ arbitrarily chosen vertices $X \subset F_1$,
        \emph{i.e.}, $E(H) \leftarrow E(H) \cup \{ (\ell, x) \mid x \in X \}$.
        \item For each vertex $v \in V(H) \setminus X \eqqcolon Y$, introduce a new vertex $w_v$ and add edge $\{v, w_v\}$, \emph{i.e.},
              $V(H) \leftarrow V(H) \cup \{ w_v \mid v \in Y \}$ and $E(H) \leftarrow E(H) \cup \{ (v, w_v) \mid v \in Y \}$,
              mark the set of $w_v$ for each $v \in Y$ as $F_2$.
    \end{enumerate}
    An example of such construction can be seen in Figure \ref{fig:proofDB}.

    Note that step 4 can always be done because $r \ge k$, so $|F_1| = n > n - r + (k - 1)$
    and since $\deg(\ell) = 0$ (before step 4), there are enough vertices for $\ell$ to connect with.
    Considering $G$ is $r$-regular, its complement, constructed and added into $H$ in steps 2 and 3, must be $(n-r-1)$-regular.
    After connecting all the $F_1$ vertices either with $\ell$, or the corresponding $w_v$ in steps 4 and 5,
    they all end up with degree $n-r$; in other words, graph $H[F_1]$ is $(n-r)$-regular.
    Also, note that the construction of $H$ is done in a time polynomial with $n$.

    Next, notice that the degree of the only leader $\ell$ is $\lambda$, as presented above.
    The other vertices $V(H) - \ell$ are naturally followers of which we can recognize two types, $F_1$ and $F_2$,
    where $\forall_{f_1 \in F_1} \deg(f_1)$ = $n-r$ and $\forall_{f_2 \in F_2} \deg(f_2)$ = $1$.
    Whereas $F_1$ are vertices of the original graph~$G$,
    $F_2$ plays the role of ``partners'' of vertices $Y$, since their only job is to substitute a missing connection with $\ell$.

    Because there is just one leader in $\mathcal{I}$, it is clear that for any follower $f' \in F'$ applies that
    $\deg(f') \geq \lambda = n - r + (k - 1)$.
    Also, because $n > r$, $\max\limits_{\forall f_1 \in F_1}\deg(f_1) = n-r > 1 = \max\limits_{\forall f_2 \in F_2}\deg(f_2)$.
    
    We now show how finding a $k$-clique in graph $G$ corresponds to finding a solution $W$ of \HLshort in~$H$, more precisely,
    we show that $(G, k)$ is a yes-instance of the $k$-\textsc{Clique} problem if and only if $\mathcal{I}$ is a yes-instance of \HLdeg.

    Let us start with the left-to-right direction, assuming that $(G, k)$ is a $yes$-instance of $k$-\textsc{Clique}.
    We denote a set of vertices that form the clique in $G$ as $S$, $|S| = k$.
    We can see from the construction above that $S \subseteq F_1$ in $H$ and also, because $S$ form a clique in $G$,
    there will be no edge connecting two vertices from $S$ in $H$.
    We now make $H[S]$ a complete graph by adding a set of new edges $E_S$, $|E_S| = \frac{k\cdot(k-1)}{2} = b$,
    which rise the degree of every follower from $S$ from $n-r$ to $n-r+(k-1) = \lambda$.
    Because $S \subseteq F_1$, $|S| = k$ and $\forall s \in S \colon \deg(s) = \lambda$,
    we get that $E_S$ is the solution $W$ of the \HLshort instance $\mathcal{I}$ and $S$ is the corresponding $F'$. 

    Next, we show the right-to-left direction of the proof, for which we assume that $\mathcal{I}$ is a yes-instance of \HLdeg,
    with newly added edges $W$ and the corresponding set of followers $F'$ as the solution. 
    For that, let us first prove three useful lemmas.

    % lemma 1
    \begin{lemma}\label{lemmaInProof}
        $F'$ can contain no vertex from $F_2$, i.e., $F' \cap F_2 = \emptyset$.
    \end{lemma}
    \begin{subproof}
        Assume, for sake of contradiction, that $F' \cap F_2 = \{f_2$\}.
        
        Because $|F'| \geq k$, there is at least $k-1$ vertices from $F_1$ in $F'$.
        Every vertex from $F_1$ have degree $n-r$ and every vertex from $F_2$ have degree $1$, so,
        because degree of each vertex from $F'$ is at least $\lambda$,
        each vertex from $F' \cap F_1$ must have gotten at least $\lambda - (n - r) = k - 1$ new neighbors and
        vertex $f_2$ must have gotten at least $\lambda - 1 = n - r + k - 2$ new neighbors,
        for which we must have added at least
        $\frac{(k - 1)^2 + (n -r + k - 2)}{2} = b + \frac{n - r - 1}{2}$ new edges into $H$.
        Because $n-1 \geq r$, we see that $\frac{n - r - 1}{2} > 0$,
        so we must have added more edges than there are in the budget, which is a contradiction.
    \end{subproof}

    We have shown that we can only add edges between followers $F_1$ when looking for solution~$W$.
    We now show that there are exactly $k$ followers in $F'$ and that
    once we add a new edge between two followers, they must become part of $F'$.

    % lemma 2
    \begin{lemma}
        There are exactly $k$ followers in $F'$, i.e., $|F'| = k$,
    \end{lemma}
    \begin{subproof}
        There cannot be less than $k$ followers in $F'$ because by the definition,
        there are at least $d = k$ followers in $F'$, so, assume, for sake of contradiction, that $|F'| > k$.

        Each vertex from $F'$ must have gotten at least $k - 1$ new neighbors,
        for which we must have added at least $\frac{|F'| \cdot (k-1)}{2} > \frac{k \cdot (k-1)}{2} = b$,
        which is a contradiction.
    \end{subproof}

    % lemma 3
    \begin{lemma}\label{lemmaInProof2}
        $W$ can contain no edge incident to a vertex that is not in $F'$, i.e.,
        $\forall (u,v) \in W \colon u \notin F_1 \setminus F' \wedge v \notin F_1 \setminus F'$.
    \end{lemma}
    \begin{subproof}
        Assume, for sake of contradiction, that $(f',f_1) \in W$, where $f' \in F'$ and $f_1 \in F_1 \setminus F'$.
        Also assume that $(f',f_1)$ is the only edge from $W$ which is incident to someone from $F_1 \setminus F'$.
    
        We know that each of the $k-1$ vertices from $F' \setminus \{f'\}$ has only gotten new neighbors from $F'$.
        Because each vertex from $F'$ must have gotten at least $k-1$ new neighbors,
        all of the vertices from $F' \setminus \{f'\}$ must have become neighbors with all the other vertices from $F'$,
        which means that $\frac{k \cdot (k-1)}{2} = b$ new edges have been added.
        However, $(f',f_1)$ is also in $W$, meaning that there are $b+1$ edges in $W$, which is a contradiction.
    \end{subproof}

    Putting it together, we can see that there are exactly $k$ followers from $F_1$ in $F'$
    and that there are only edges between followers from $F'$ in $W$.
    
    To get at least $k-1$ new neighbors, each follower from $F'$ has connected with all the other
    followers from $F'$, meaning that $H[F']$ is a complete graph.
    Also note that every edge from $W$ must be an edge from the original graph $G$ because $W \subset F_1 \times F_1$ and $F_1 = V(G)$
    and also all the edges from $E(\overline{G})$ are already in $E(H)$ so they cannot be in $W$, since it only contains newly added edges.
    For this reason, the fact that $H[F']$ is a complete graph means that vertices $F'$ form a $k$-clique in~$G$.

    Because a solution $W$ for \HLshort in $H$ exists if and only if a $k$-clique in $G$ exists,
    we can conclude that finding a solution for the \HL problem in graph $H$ is at least as hard as
    finding a solution for the $k$-\textsc{Clique} problem in graph $G$.

    The reduction presented in this proof is a valid parameterized reduction because;
    $(G, k)$ is a yes-instance of the $k$-\textsc{Clique} problem if and only if $\mathcal{I}$ is a yes-instance of \HLdeg,
    the construction of $H$ is done in polynomial time;
    and there is a function of $k$, $g(k) = \frac{k\cdot(k-1)}{2} + k$, upper-bounding the parameter $k'$ of $\mathcal{I}$.
\end{proof}

% Figure - proof 1
\begin{figure}[t!]
    \centering
    % G
    \begin{subfigure}[t]{.46\textwidth}
        \centering
        \begin{tikzpicture}[node distance={16mm}, scale=0.8, every node/.style={scale=0.8}]
            \tikzstyle{vertex} = [circle, draw=black, text=white]
            \tikzstyle{edge} = []
            
            \node[vertex] (g1) at (0,0) {$f_{1_1}$};
            \node[vertex] (g2) [right of = g1] {$f_{1_2}$};
            \node[vertex] (g3) [below of = g1] {$f_{1_3}$};
            \node[vertex] (g4) [right of = g3] {$f_{1_4}$};
            \node[vertex] (g5) [below of = g3] {$f_{1_5}$};
            \node[vertex] (g6) [right of = g5] {$f_{1_6}$};
            \draw[edge] (g1)--(g2);
            \draw[edge] (g5)--(g6);
            \draw[edge] (g1)--(g3);
            \draw[edge] (g3)--(g5);
            \draw[edge] (g2)--(g4);
            \draw[edge] (g4)--(g6);
            \draw[edge] (g3)--(g2);
            \draw[edge] (g3)--(g6);
            \draw[edge] (g4)--(g1);
            \draw[edge] (g4)--(g5);
            \draw[edge] (g1) to [out=180,in=180,looseness=0.5] (g5);
            \draw[edge] (g2) to [out=0,in=0,looseness=0.5] (g6);
            % dummy to ensure the same height as the other subfig
            \node (l) [right of = g4] {};
            \draw[color=white] (l) to [out=315,in=270] (g5);
        \end{tikzpicture}
        \caption{
            ~A $4$-regular ($r = 4$) graph $G$ with $n=6$ vertices,
            playing a role of the input graph for the $k$-\textsc{Clique} problem and
            is the starting point of our construction.
        }
    \end{subfigure}%
    ~~~~
    % H
    \begin{subfigure}[t]{.46\textwidth}
        \centering
        \begin{tikzpicture}[node distance={16mm}, scale=0.8, every node/.style={scale=0.8}]
            \tikzstyle{followerG} = [circle, draw=black]
            \tikzstyle{edgeG} = [dashed, draw=green!40, thick]

            \tikzstyle{edgeCompl} = []

            \tikzstyle{followerF2} = [circle, fill=red!50]
            \tikzstyle{edgeF2} = [draw=red!50]

            \tikzstyle{square} = [regular polygon, regular polygon sides=4]
            \tikzstyle{leader} = [square, fill=blue!50, label={[text=blue!50]above:$\ell$}]
            \tikzstyle{edgel} = [color=blue!50]

            % G in H
            \node[followerG] (g1) at (0,0) {$f_{1_1}$};
            \node[followerG] (g2) [right of = g1] {$f_{1_2}$};
            \node[followerG] (g3) [below of = g1] {$f_{1_3}$};
            \node[followerG] (g4) [right of = g3] {$f_{1_4}$};
            \node[followerG] (g5) [below of = g3] {$f_{1_5}$};
            \node[followerG] (g6) [right of = g5] {$f_{1_6}$};
            \draw[edgeG] (g1)--(g2);
            \draw[edgeG] (g5)--(g6);
            \draw[edgeG] (g1)--(g3);
            \draw[edgeG] (g3)--(g5);
            \draw[edgeG] (g2)--(g4);
            \draw[edgeG] (g4)--(g6);
            \draw[edgeG] (g3)--(g2);
            \draw[edgeG] (g3)--(g6);
            \draw[edgeG] (g4)--(g1);
            \draw[edgeG] (g4)--(g5);
            \draw[edgeG] (g1) to [out=180,in=180,looseness=0.5] (g5);
            \draw[edgeG] (g2) to [out=0,in=0,looseness=0.5] (g6);
            % complement
            \draw[edgeCompl] (g1)--(g6);
            \draw[edgeCompl] (g2)--(g5);
            \draw[edgeCompl] (g3)--(g4);
            % leader
            \node[leader] (l) [right of = g4] {};
            \draw[edgel] (l)--(g4);
            \draw[edgel] (l)--(g6);
            \draw[edgel] (l) to [out=45,in=90] (g1);
            \draw[edgel] (l) to [out=315,in=270] (g5);
            % partners
            \node[followerF2, label={[text=red!50]above:$f_{2_1}$}] (f2) [above of = g2] {};
            \node[followerF2, label={[text=red!50]above:$f_{2_2}$}] (f3) [left of = g3] {};
            \draw[edgeF2] (f2)--(g2);
            \draw[edgeF2] (f3)--(g3);
        \end{tikzpicture}
        \caption{
            ~A graph $H$ constructed from graph $G$ where $H[F_1]$ is $2$-regular.
            The black vertices are followers from $F_1$,
            the red vertices are followers from $F_2$ and
            the blue vertex is a leader.
            The black edges are edges $E(\overline{G})$ and the green, dashed lines represent edges from $G$,
            which are not present in $H$, but from which a potential solution $W$ of \HLshort would be picked.
        }
    \end{subfigure}
    \caption{A sample construction of graph $H$ as presented in Proof \ref{proofDB}.}
    \label{fig:proofDB}
\end{figure}


% Corollaries
We have just shown that \HL parameterized by $b+d$ is \Wh even with a constant number of leaders.
We can easily see that \HLshort remains \Wh even when parameterized by $|L|+b+d$.
This is because the problem is \Wh already for $|L|$ bounded by a constant function,
hence, bounding it by a nonconstant function will not make the problem easier to solve, leaving it \Wh.

\begin{remark}\label{cor:LnBD:Wh}
    \HL parameterized by $|L|+b+d$ is \Wh.
\end{remark}

Following the argument, we can see that \HL is \Wh for any variation of parameters $|L|$, $b$, $d$.
For example, if we omit $b$ from the parameter, then $b$ is no longer bounded by a function of the parameter and
become bounded by a function of the input size, making the problem possibly harder and so leaving it \Wh.

\begin{corollary}\label{cor:LnBD:variation:Wh}
    \HL parameterized by any variation of $|L|$, $b$, $d$ is \Wh.
\end{corollary}

The parameterized reduction shown in Proof \ref{proofDB} is done in polynomial time with respect to input size,
that means that it is also a polynomial reduction.
That is, there is a polynomial reduction from $k$-\textsc{Clique} on regular graphs to \HL, where $|L| = 1$.
Because $k$-\textsc{Clique} on regular graphs is \NPh, we get from the reduction that \HL is \NPh as well,
even if the number of leaders is constant.
From this, we get the following.

\begin{corollary}\label{cor:Ln:pNPh}
    \HL is \pNPh with respect to $|L|$.
\end{corollary}


% XP algo for param b
We now look at the problem from the opposite end.
Let us recall that we interpret \W-hardness as an evidence that the given problem is not in \FPT.
Hence, there is no use in trying to find FPT algorithms for \HL parameterized by $|L|$, $b$, $d$, or any of their variation.
However, there is still a chance for XP algorithms.
We now show that there are indeed such algorithms for parameters $b$ and $d$, meaning that
\HL parameterized either by budget $b$ or by safety margin $d$ belongs to \XP.
From this, we then show that \HLshort is in \XP even for the other parameter variations.
Let us present a proof of the variant with parameter $b$ first.

\begin{theorem}\label{theorem:B:XP}
    \HL parameterized by $b$ is in \XP.
\end{theorem}

\begin{proof}
    Given parameterized problem $(\mathcal{I}, b)$, where $\mathcal{I} = (G, L, b, c_{deg}, d)$ is a \HLdeg instance
    and~$b$ is the parameter,
    we show an algorithm which decides whether or not $(\mathcal{I}, b)$ is a $yes$-instance in time bounded by
    $f(b) \cdot n^{g(b)}$ for some computable functions $f,g$.

    We start by picking gradually each of the $\binom{|\hat{A}|}{b}$,
    $\binom{|\hat{A}|}{b} \leq \binom{n^2}{b} \leq \frac{n^{2b}}{b!} \leq n^{\mathcal{O}(b)}$,
    $b$-element subsets of edges that can be added between followers, we denote the current subset as $B$.
    Because we take all the possible $b$-element subsets of $\hat{A}$,
    we eventually find a subset $B$ such that $B = W$;
    if $(\mathcal{I}, b)$ is a $yes$-instance, then such $W$ exists.
    We then try to add those edges $B$ into $G$ and after each such addition,
    we check whether there are at least $d$ followers
    with degree greater than~$\lambda$.
    The addition can be done in time $\mathcal{O}(b)$, and the check can be done in time $\mathcal{O}(n)$.

    If instance $(\mathcal{I}, b)$ is not identified as a $yes$-instance for any choice of $B$,
    then we conclude that $(\mathcal{I}, b)$ is a $no$-instance because if it was a $yes$-instance,
    then we would have found it by the procedure above.

    The algorithm we have just described runs in time $n^{\mathcal{O}(b)}$.
    From this, we can conclude that \HL parameterized by $b$ is in \XP.
\end{proof}

Notice that although the just presented algorithm is not possibly the best one we could find,
there surely cannot be so much faster algorithm since, as we already know,
there can be no FPT algorithm for \HL parameterized by $b$.


% XP algo for param d
Now we proceed to a proof of the variant with parameter $d$.
This proof will be led in a similar manner as the previous proof.

\begin{theorem}\label{theorem:D:XP}
    \HL parameterized by $d$ is in \XP.
\end{theorem}

\begin{proof}\label{proof:XPd}
    Given parameterized instance $(\mathcal{I}, d)$, where $\mathcal{I} = (G, L, b, c_{deg}, d)$ is a \HLdeg instance and~$d$ is the parameter,
    we show an algorithm which decides whether or not $(\mathcal{I}, d)$ is a $yes$-instance in time bounded by
    $f(d) \cdot n^{g(d)}$ for some computable functions $f,g$.

    We start by picking gradually each of the $\binom{|F|}{d}$,
    $\binom{|F|}{d} \leq \binom{n}{d} \leq \frac{n^{d}}{d!} \leq n^{\mathcal{O}(d)}$,
    $d$-element subsets of followers $F$,
    we denote the current subset as $D$ and a set of edges that can be added between followers from $D$
    as $\hat{A}_D \subseteq \hat{A}$, $\hat{A}_D = E(\overline{G[D]}) \leq d^2$.
    Because we take all the possible $d$-element subsets of $F$, we eventually find a subset $D$ such that,
    after the addition of new edges, $D = F'$;
    if such $F'$ exists, that is, if $(\mathcal{I}, d)$ is a $yes$-instance.
    We then try to add $t \eqqcolon \min(|\hat{A}_D|, b)$ edges between followers from $D$
    in each of the $\binom{|\hat{A}_D|}{t} \leq 2^{d^2}$ possible ways.
    Note that $t$ can be equal to $b$ because we are not allowed to add more than $b$ edges into $G$.
    The addition can be done in time $\mathcal{O}(t)$.
    From now, only one of these three situations can occur:

    First, we added $t$ edges into $G$ and all $D$ vertices have a degree greater than $\lambda$.
    Because $|D| = d$, $(\mathcal{I}, d)$ is a $yes$-instance.
    We can check this in time $\mathcal{O}(n)$.
    
    Second, we added $t$ edges into $G$ but not all $D$ vertices have a degree greater than $\lambda$.
    Moreover, $|\hat{A}_D| < b$ and thus we know that $G[D]$ is (after adding new edges) already a complete graph
    but we still can add some edges into $G$ before reaching the budget $b$. 
    In this situation, $D = F'$ if and only if there is at least $\lambda - \min\limits_{\forall \delta \in D}deg(\delta)$ followers
    outside of $D$ and a number of edges we are still allowed to add before reaching budget $b$, $b - |\hat{A}_D|$,
    is smaller than the sum of the number of connections each vertex from $D$ needs to reach $\lambda$,
    $\sum_{\delta \in D} \min(0, \lambda - \deg(\delta))$.
    This is true because in order for instance $(\mathcal{I}, d)$ to be a $yes$-instance, we must
    ensure that each follower from $D$ has degree at least $\lambda$, for what we connect them to some followers
    outside $D$ (because $D$ is already complete), for which there must be enough followers outside $D$ and
    we must have enough edges left before reaching $b$.
    Note that if $D = F'$, then $(\mathcal{I}, d)$ in a $yes$-instance.
    We can check this in time $\mathcal{O}(n)$.
    
    Third, we added $t$ edges into $G$ but not all $D$ vertices have a degree greater than $\lambda$.
    Moreover, $|\hat{A}_D| \geq b$, so $t = b$ and thus the current set $D$ can never be $F'$.
    We can check this in time $\mathcal{O}(n)$.

    Situations other than one of the three cannot occur.
    If instance $(\mathcal{I}, d)$ is not identified as a $yes$-instance for any choice of $D$,
    then we conclude that $(\mathcal{I}, d)$ is a $no$-instance because if it was a $yes$-instance,
    then we would have found it by the procedure above.
    
    The algorithm we have just described runs in time $\mathcal{O}(2^{d^2}) \cdot n^{\mathcal{O}(d)}$.
    From this, we can conclude that \HL parameterized by $d$ is in \XP.
\end{proof}

We have just shown that \HL parameterized by $b$ or $d$ is in \XP.
If we now add some other parameter or parameters to either $b$ or $d$, then the problem
might become easier but surely stays in \XP.
For example, take $b+|L|$ as the parameter;
when not a part of the parameter, |L| was bounded by a function of input size,
because it is now a part of the parameter, |L| is bounded by a function of the parameter,
hence the problem is possibly easier to solve.
From this, we get the following.

\begin{corollary}\label{cor:BD:variation:XP}
    \HL is in \XP when parameterized by at least $b$ or $d$.
\end{corollary}

Last, following the same logic and considering the results from Dey and Medya \cite{Dey2019} presented in Subsection \ref{subsection:ResultsDegree},
we can state the following.

\begin{corollary}\label{cor:Ld:variation:FPT}
    \HL is in \XP when parameterized by at least $\lambda$.
\end{corollary}

%---------------------------------------------------------------
\section{Overview}

In this last section, we visualize the results for \HLdeg presented in Section \ref{section:OurResults},
together with the result for parameter $\lambda$ presented in Subsection \ref{subsection:ResultsDegree}.
The visualization can be seen in Figure \ref{fig:complexityPicture}.

% Figure - overview
\begin{figure}[t]
    \centering
    \begin{tikzpicture}[node distance=18mm]
        \tikzstyle{every node} = [minimum height=0.7cm,minimum width=0.9cm,scale=1.3];
        \tikzstyle{result} = [draw, very thick];
        \tikzstyle{corollary} = [draw,dashed];
        \tikzstyle{pNPh} = [fill=red!30];
        \tikzstyle{Wh} = [fill=orange!30];
        \tikzstyle{XP} = [fill=green!30];

        \tikzstyle{edgeXP} = [color=orange,-{Stealth[length=3mm]}]
        \tikzstyle{edgeW} = [color=violet,-{Stealth[length=3mm]}]
        \tikzstyle{edgeFPT} = [color=green,-{Stealth[length=3mm]}]

        % pairs
        \node[corollary,Wh,label={0:\small{\tt C\ref{cor:LnBD:variation:Wh}}}] (LnB) at (0,0) {\leadnum, \budget};
        \node[corollary,Wh,label={0:\small{\tt C\ref{cor:LnBD:variation:Wh}}}] (LnD) [right =of LnB] {\leadnum, \safetymarg};
        \node[corollary,Wh,label={0:\small{\tt C\ref{cor:LnBD:variation:Wh}}}] (BD) [right=of LnD] {\budget, \safetymarg};
        % singletons
        \node[result,Wh,label={0:\small{\tt T\ref{theorem:B:XP}}}] (B) [above left=18mm and 2mm of LnD] {\budget};
        \node[result,Wh,label={0:\small{\tt T\ref{theorem:D:XP}}}] (D) [above right=18mm and 2mm of LnD] {\safetymarg};
        \node[corollary,pNPh,label={0:\small{\tt C\ref{cor:Ln:pNPh}}}] (Ln) [left=of B] {\leadnum};
        \node[XP,label={0:\small{\tt known\cite{Dey2019}}}] (Ld) [right=of D] {\leaddeg};
        % triplets
        \node[result,Wh,label={0:\small{\tt T\ref{theorem:DB}}}] (LnBD) [below=of LnD] {\leadnum, \budget, \safetymarg};
        % edges
        \draw[->,edgeXP] (B.225) -- (LnB.135);
        \draw[->,edgeXP] (D.225) -- (LnD.135);
        \draw[->,edgeXP] (D.225) -- (BD.135);    
        \draw[->,edgeXP] (BD.225) -- (LnBD.135);
        
        \draw[->,edgeW] (LnBD.45) -- (LnB.315);
        \draw[->,edgeW] (LnBD.45) -- (LnD.315);
        \draw[->,edgeW] (LnBD.45) -- (BD.315);
        \draw[->,edgeW] (BD.45) -- (B.315);
        \draw[->,edgeW] (BD.45) -- (D.315);
    \end{tikzpicture}
    \caption
    [Overviewing graph with results for \HLdeg]
    {
        The complexity picture for the \HL problem given the degree centrality measure, \HLdeg.
        Combinations of parameters that give rise to XP algorithms are highlighted in green,
        while combinations for which \HLdeg is \Wh but still in \XP are highlighted in orange and \pNP combinations
        are highlighted in red.
        Results explicitly proved in this work are represented by a black solid border,
        their corollaries are represented by a black dashed border and
        results from the literature do not have any border.
        Each node has a reference to the corresponding theorem, remark, corollary or bibliography record.
        Arrows represent how complexity results spread.
        Orange arrows represent \XP class and violet arrows represent \W-hardness.
        When multiple arrows of the same color would aim at the same node,
        then only one of them is displayed.  
    }
    \label{fig:complexityPicture}
\end{figure}
%---------------------------------------------------------------

%---------------------------------------------------------------
\chapter{Conclusion and future work}
%---------------------------------------------------------------

In this thesis, we introduced the reader to the topic of covert networks and their analysis,
with emphasis on the \HL problem for which we provided a detailed description, we motivated its study
and explained its hardness.
We mostly followed up on the work from Waniek et al. and Dey with Medya \cite{Waniek2017,Dey2019}.
We pointed out the authors use definitions that are different from each other and described our motivation
for picking the one from Dey and Medya.
We then reviewed the result for the degree centrality from Waniek et al. with respect to our definition.

The results from the literature are described in a language of classical complexity theory,
so we put them into a perspective of the parameterized complexity framework, which we also briefly introduced.
We then focused on the problem with respect to the degree centrality measure
and presented our own parameterized complexity results in this domain.
Namely, we showed that \HL is \Wh when parameterized by $b+d$ even if $|L| = 1$.
Because the reduction in the proof was done in polynomial time,
we got that \HLshort is \pNPh with respect to $|L|$.
We then showed that \HLshort is in \XP when parameterized by $b$ or $d$.
Last, we put together both our results and results from the literature for the degree centrality
and parameters $|L|$, $b$, $d$, $\lambda$ and visualized them in an overviewing graph.

Some of the future work on the topic of the \HL problem may involve
revisioning the results for the closeness and betweenness centralities from Waniek et al. \cite{Waniek2021full}
with respect to our definition,
or inspecting the problem for various combinations of centralities and parameters for which there are no results yet.

For instance, it would be interesting to see some parameterized complexity results for the core centrality
as the problem seems much harder for the core centrality than for the degree centrality \cite{Dey2019}.
For the variant with the core centrality, one could also survey the computational complexity of \HLshort
when the core centrality of every leader is at most 2,
as Dey and Medya \cite{Dey2019} inspect only the situation where the core centrality of every leader is exactly 3.

For the variant with the degree centrality, which we studied the most in this thesis,
it would be interesting to see the parameterization which gives rise to an FPT algorithm
as there is known no such algorithm yet.
It seems that the parameterization by the \emph{vertex cover number} could, in some cases,
bring promising results in this direction.
Vertex cover number is the size of a \emph{minimum vertex cover},
minimum vertex cover is the smallest possible number of vertices $V$ of a given graph $G$ such that each edge from $E(G)$ is incident
to at least one vertex from $V$.
For the parameter $\lambda$, \HLdeg is in \XP and
we think that it is also \Wh, so, another immediate future work is to prove or disprove this conjecture.

The \HL problem has not yet been studied parameterized by some structural limitation of the input graph.
One such structural limitation is the aforementioned \emph{minimum vertex cover}.
Another interesting structural parameter to expose the problem to is a \emph{treewidth}.
Treewidth is, roughly speaking, a measurement of how well a graph can be decomposed into pieces that are connected in a
tree-like fashion \cite[p.~151]{Cygan2015}.
The notion of treewidth was described by Robertson and Seymour in the series of papers \cite{Robertson1984,Robertson1986.2,Robertson1986.5}. 
Treewidth has vast applications in the design of parameterized graph algorithms and
in the design of graph algorithms in general.

Another open problem in this topic is exploring the average case computational complexity of \HL for various centrality measures,
rather than doing the worst-case analysis.
The results from the literature show that the problem is intractable only in the worst case so
some heuristics may exist that efficiently solve most of the instances typical in some real-world application.
If true, then the apparent complexity of hiding in networks by manipulating their structure
could be overcome in many real-world cases.~\cite{Dey2019}

%\include{text/samples/cviko.tex}
%% Do not forget to include Introduction
%---------------------------------------------------------------
% \chapter{Introduction}
% uncomment the following line to create an unnumbered chapter
\chapter*{Introduction}\addcontentsline{toc}{chapter}{Introduction}\markboth{Introduction}{Introduction}
%---------------------------------------------------------------
\setcounter{page}{1}

% The following environment can be used as a mini-introduction for a chapter. Use that any way it pleases you (or comment it out). It can contain, for instance, a summary of the chapter. Or, there can be a quotation.
\begin{chapterabstract}
	\lipsum[1]
\end{chapterabstract}

\lipsum[2][1-4]{} [1]

\lipsum[4]

%---------------------------------------------------------------
\section{Ut enim ad minim veniam}
%---------------------------------------------------------------

\lipsum[6-7]

\begin{figure}
\centering
%\includegraphics[scale=0.4]{pic/index}
\resizebox{\textwidth}{!}{
\begin{tikzpicture}[level/.style={sibling distance=60mm/#1}]
\node [circle,draw] (z){$n$}
  child {node [circle,draw] (a) {$\frac{n}{2}$}
    child {node [circle,draw] (b) {$\frac{n}{2^2}$}
      child {node {$\vdots$}
        child {node [circle,draw] (d) {$\frac{n}{2^k}$}}
        child {node [circle,draw] (e) {$\frac{n}{2^k}$}}
      } 
      child {node {$\vdots$}}
    }
    child {node [circle,draw] (g) {$\frac{n}{2^2}$}
      child {node {$\vdots$}}
      child {node {$\vdots$}}
    }
  }
  child {node [circle,draw] (j) {$\frac{n}{2}$}
    child {node [circle,draw] (k) {$\frac{n}{2^2}$}
      child {node {$\vdots$}}
      child {node {$\vdots$}}
    }
  child {node [circle,draw] (l) {$\frac{n}{2^2}$}
    child {node {$\vdots$}}
    child {node (c){$\vdots$}
      child {node [circle,draw] (o) {$\frac{n}{2^k}$}}
      child {node [circle,draw] (p) {$\frac{n}{2^k}$}
        child [grow=right] {node (q) {$=$} edge from parent[draw=none]
          child [grow=right] {node (q) {$O_{k = \lg n}(n)$} edge from parent[draw=none]
            child [grow=up] {node (r) {$\vdots$} edge from parent[draw=none]
              child [grow=up] {node (s) {$O_2(n)$} edge from parent[draw=none]
                child [grow=up] {node (t) {$O_1(n)$} edge from parent[draw=none]
                  child [grow=up] {node (u) {$O_0(n)$} edge from parent[draw=none]}
                }
              }
            }
            child [grow=down] {node (v) {$O(n \cdot \lg n)$}edge from parent[draw=none]}
          }
        }
      }
    }
  }
};
\path (a) -- (j) node [midway] {+};
\path (b) -- (g) node [midway] {+};
\path (k) -- (l) node [midway] {+};
\path (k) -- (g) node [midway] {+};
\path (d) -- (e) node [midway] {+};
\path (o) -- (p) node [midway] {+};
\path (o) -- (e) node (x) [midway] {$\cdots$}
  child [grow=down] {
    node (y) {$O\left(\displaystyle\sum_{i = 0}^k 2^i \cdot \frac{n}{2^i}\right)$}
    edge from parent[draw=none]
  };
\path (q) -- (r) node [midway] {+};
\path (s) -- (r) node [midway] {+};
\path (s) -- (t) node [midway] {+};
\path (s) -- (l) node [midway] {=};
\path (t) -- (u) node [midway] {+};
\path (z) -- (u) node [midway] {=};
\path (j) -- (t) node [midway] {=};
\path (y) -- (x) node [midway] {$\Downarrow$};
\path (v) -- (y)
  node (w) [midway] {$O\left(\displaystyle\sum_{i = 0}^k n\right) = O(k \cdot n)$};
\path (q) -- (v) node [midway] {=};
\path (e) -- (x) node [midway] {+};
\path (o) -- (x) node [midway] {+};
\path (y) -- (w) node [midway] {$=$};
\path (v) -- (w) node [midway] {$\Leftrightarrow$};
\path (r) -- (c) node [midway] {$\cdots$};
\end{tikzpicture}}
\caption{~Lorem ipsum dolor sit amet}\label{img:index}
\end{figure}

%---------------------------------------------------------------
\section{Ut enim ad minim veniam}
%---------------------------------------------------------------

\lipsum[2-4]

%---------------------------------------------------------------
\subsection{Ut enim ad minim veniam}
%---------------------------------------------------------------

Curabitur ligula sapien, pulvinar a vestibulum quis, facilisis vel sapien. Duis condimentum augue id magna semper rutrum. Aliquam ornare wisi eu metus. Fusce aliquam vestibulum ipsum. Vivamus ac leo pretium faucibus\ref{img:index}.

\begin{itemize}
    \item Ut enim ad minim veniam, quis nostrud
    \item Ut enim ad minim 
    \item Ut enim ad minim veniam, quis 
    \begin{itemize}
        \item Ut enim ad
        \item Ut enim ad
        \begin{itemize}
            \item Ut enim 
            \item Ut enim 
            \begin{itemize}
            \item Ut enim 
            \item Ut enim 
        \end{itemize}
        \end{itemize}
    \end{itemize}
\end{itemize}

\section{Class aptent taciti}

\lipsum[2]

\subsection{Class aptent taciti}

\lipsum[6-7]

\begin{enumerate}
    \item Ut enim ad minim veniam, quis nostrud
    \item Ut enim ad minim 
    \item Ut enim ad minim veniam, quis 
    \begin{enumerate}
        \item Ut enim ad
        \item Ut enim ad
        \begin{enumerate}
            \item Ut enim 
            \item Ut enim 
            \begin{enumerate}
            \item Ut enim 
            \item Ut enim 
        \end{enumerate}
        \end{enumerate}
    \end{enumerate}
\end{enumerate}


%---------------------------------------------------------------
\section{Ut enim ad minim veniam, quis nostrud}
%---------------------------------------------------------------

Ut enim ad minim veniam, quis nostrud exercitation ullamco laboris nisi ut aliquip ex ea commodo consequat. Nulla non arcu lacinia neque faucibus fringilla. Vestibulum erat nulla, ullamcorper nec, rutrum non, nonummy ac, erat. Aliquam erat volutpat. Proin pede metus, vulputate nec, fermentum fringilla, vehicula vitae, justo.\footnote{Ut enim ad minim veniam, quis nostrud exercitation.} Etiam dictum tincidunt diam. In laoreet, magna id viverra tincidunt, sem odio bibendum justo, vel imperdiet sapien wisi sed libero. Nulla est. Maecenas fermentum, sem in pharetra pellentesque, velit turpis volutpat ante, in pharetra metus odio a lectus. Duis aute irure dolor in reprehenderit in voluptate velit esse cillum dolore eu fugiat nulla pariatur. 

\begin{lstlisting}[caption={~Zbytečný kód},label=list:8-6,captionpos=t,float,abovecaptionskip=-\medskipamount,belowcaptionskip=\medskipamount,language=C]
    #include<stdio.h>
    #include<iostream>
    // A comment
    int main(void)
    {
        printf("Hello World\n");
        return 0;
    }
\end{lstlisting}

%%%%%%%%%%%%%%%%%%%%%%%%%%%%%%%%%
% alternative using package minted for source highlighting
% package minted requires execution with `-shell-escape'
% e.g., `xelatex -shell-escape ctufit-thesis.tex'
% \begin{listing}
% \caption{Zbytečný kód}\label{list:8-6}
% \begin{minted}{C}
%     #include<stdio.h>
%     #include<iostream>
%     // A comment
%     int main(void)
%     {
%         printf("Hello World\n");
%         return 0;
%     }
% \end{minted}
% \end{listing}
% %%%%%%%%%%%%%%%%%%%%%%%%%%%%%%%%%
Nullam feugiat, turpis at pulvinar vulputate, erat libero tristique tellus, nec bibendum odio risus sit amet ante. Aenean id metus id velit ullamcorper pulvinar. Fusce wisi. Integer lacinia. Aliquam id dolor. Pellentesque pretium lectus id turpis. Suspendisse sagittis ultrices augue. In laoreet, magna id viverra tincidunt, sem odio bibendum justo, vel imperdiet sapien wisi sed libero. Sed ac dolor sit amet purus malesuada congue. \cite{Crochemore2002}

Class aptent taciti sociosqu ad litora torquent per conubia nostra, per inceptos hymenaeos. Fusce suscipit libero eget elit. Etiam dui sem, fermentum vitae, sagittis id, malesuada in, quam. Aliquam id dolor. Curabitur bibendum justo non orci. Duis viverra diam non justo. Curabitur ligula sapien, pulvinar a vestibulum quis, facilisis vel sapien. Duis condimentum augue id magna semper rutrum. Aliquam ornare wisi eu metus. Fusce aliquam vestibulum ipsum. Vivamus ac leo pretium faucibus. \cite{Motwani2014}

%---------------------------------------------------------------
\subsection{Ut enim ad minim veniam, quis nostrud}
%---------------------------------------------------------------

Ut enim ad minim veniam, quis nostrud exercitation ullamco laboris nisi ut aliquip ex ea commodo consequat. Nulla non arcu lacinia neque faucibus fringilla. Vestibulum erat nulla, ullamcorper nec, rutrum non, nonummy ac, erat. Aliquam erat volutpat. Proin pede metus, vulputate nec, fermentum fringilla, vehicula vitae, justo. Etiam dictum tincidunt diam. In laoreet, magna id viverra tincidunt, sem odio bibendum justo. \cite{Sestakova2018} 

\begin{table}\centering
\caption[Příklad tabulky]{~Zadávání matematiky}\label{tab:matematika}
\begin{tabular}{l|l|c|c}
	Typ		& Prostředí		& \LaTeX{}ovská zkratka	& \TeX{}ovská zkratka	\tabularnewline \hline 
 	Text		& \verb|math|		& \verb|\(...\)|	& \verb|$...$|	\tabularnewline \hline
 	Displayed	& \verb|displaymath|	& \verb|\[...\]|	& \verb|$$...$$|	\tabularnewline 
\end{tabular}
\end{table}


Nulla est. Maecenas fermentum, sem in pharetra pellentesque, velit turpis volutpat ante, in pharetra metus odio a lectus. Duis aute irure dolor in reprehenderit in voluptate velit esse cillum dolore eu fugiat nulla pariatur. Nullam feugiat, turpis at pulvinar vulputate, erat libero tristique tellus, nec bibendum odio risus sit amet ante. Aenean id metus id velit ullamcorper pulvinar. 

\subsubsection{Class aptent taciti}

\begin{definition}[Optional label]
Class aptent taciti sociosqu ad litora torquent per conubia nostra, per inceptos hymenaeos. Fusce suscipit libero eget elit. Etiam dui sem, fermentum vitae, sagittis id, malesuada in, quam. Aliquam id dolor. Curabitur bibendum justo non orci.
\end{definition}

\begin{example}
Class aptent taciti sociosqu ad litora torquent per conubia nostra, per inceptos hymenaeos. Fusce suscipit libero eget elit. Etiam dui sem, fermentum vitae, sagittis id, malesuada in, quam. Aliquam id dolor. Curabitur bibendum justo non orci.
\end{example}

\begin{theorem}
Class aptent taciti sociosqu ad litora torquent per conubia nostra, per inceptos hymenaeos. Fusce suscipit libero eget elit. Etiam dui sem, fermentum vitae, sagittis id, malesuada in, quam. Aliquam id dolor. Curabitur bibendum justo non orci.
\end{theorem}

\begin{proof}
Fusce suscipit libero eget elit. Etiam dui sem, fermentum vitae, sagittis id, malesuada in, quam. Aliquam id dolor. Curabitur bibendum justo non orci.
\end{proof}

\begin{corollary}
Fusce suscipit libero eget elit. Etiam dui sem, fermentum vitae, sagittis id, malesuada in, quam. Aliquam id dolor. Curabitur bibendum justo non orci.
\end{corollary}

\begin{proposition}
Fusce suscipit libero eget elit. Etiam dui sem, fermentum vitae, sagittis id, malesuada in, quam. Aliquam id dolor. Curabitur bibendum justo non orci.
\end{proposition}

\begin{note}
Fusce suscipit libero eget elit. Etiam dui sem, fermentum vitae, sagittis id, malesuada in, quam. Aliquam id dolor. Curabitur bibendum justo non orci.
\end{note}

\begin{remark}
Fusce suscipit libero eget elit. Etiam dui sem, fermentum vitae, sagittis id, malesuada in, quam. Aliquam id dolor. Curabitur bibendum justo non orci.
\end{remark}

\begin{lemma}
Class aptent taciti sociosqu ad litora torquent per conubia nostra, per inceptos hymenaeos. Fusce suscipit libero eget elit. Etiam dui sem, fermentum vitae, sagittis id, malesuada in, quam. Aliquam id dolor. Curabitur bibendum justo non orci.
\end{lemma}

\lipsum[1-2]

\subsection{Class aptent taciti sociosqu}

\lipsum[4-5]

%---------------------------------------------------------------
\chapter{Lorem ipsum}
%---------------------------------------------------------------

\begin{chapterabstract}
	Lorem ipsum dolor sit amet, consectetuer adipiscing elit. Curabitur sagittis hendrerit ante. Class aptent taciti sociosqu ad litora torquent per conubia nostra, per inceptos hymenaeos. Cras pede libero, dapibus nec, pretium sit amet, tempor quis. Sed vel lectus. Donec odio tempus molestie, porttitor ut, iaculis quis, sem. Cras pede libero, dapibus nec, pretium sit amet, tempor quis. Sed vel lectus. 
\end{chapterabstract}

Lorem ipsum dolor sit amet, consectetuer adipiscing elit. Curabitur sagittis hendrerit ante. Class aptent taciti sociosqu ad litora torquent per conubia nostra, per inceptos hymenaeos. Cras pede libero, dapibus nec, pretium sit amet, tempor quis. Sed vel lectus. Donec odio tempus molestie, porttitor ut, iaculis quis, sem. Suspendisse sagittis ultrices augue. Donec ipsum massa, ullamcorper in, auctor et, scelerisque sed, est. In sem justo, commodo ut, suscipit at, pharetra vitae, orci. Pellentesque pretium lectus id turpis. \cite{Kopka2004}

\section{Donec odio tempus molestie}

\lipsum[2] \cite{def:1, def:2}

\subsection{Class aptent taciti}

\lipsum[2-3]

\begin{description}
\item[Kapitola 1] Lorem ipsum dolor sit amet, consectetuer adipiscing elit. Curabitur sagittis hendrerit ante. Class aptent taciti sociosqu ad litora torquent per conubia nostra, per inceptos hymenaeos. Cras pede libero, dapibus nec, pretium sit amet, tempor quis.

\item[Kapitola 2] Lorem ipsum dolor sit amet, consectetuer adipiscing elit. Curabitur sagittis hendrerit ante. Class aptent taciti sociosqu ad litora torquent per conubia nostra, per inceptos hymenaeos. Cras pede libero, dapibus nec, pretium sit amet, tempor quis.

\item[Kapitola 3] Lorem ipsum dolor sit amet, consectetuer adipiscing elit. Curabitur sagittis hendrerit ante. Class aptent taciti sociosqu ad litora torquent per conubia nostra, per inceptos hymenaeos. Cras pede libero, dapibus nec, pretium sit amet, tempor quis.

\item[Kapitola 4] Lorem ipsum dolor sit amet, consectetuer adipiscing elit. Curabitur sagittis hendrerit ante. Class aptent taciti sociosqu ad litora torquent per conubia nostra, per inceptos hymenaeos. Cras pede libero, dapibus nec, pretium sit amet, tempor quis.
\end{description}

\lipsum[2]

\section{Lorem ipsum dolor sit amet}

\lipsum[3-5]


\appendix\appendixinit % do not remove these two commands

% \chapter{Some appendix}


Sem přijde to, co nepatří do hlavní části.
 % include `appendix.tex' from `text/' subdirectory

\backmatter % do not remove this command

\printbibliography % print out the BibLaTeX-generated bibliography list

% \chapter{Obsah přiloženého média}


	\dirtree{%
		.1 readme.txt\DTcomment{stručný popis obsahu média}.
		.1 exe\DTcomment{adresář se spustitelnou formou implementace}.
		.1 src.
		.2 impl\DTcomment{zdrojové kódy implementace}.
		.2 thesis\DTcomment{zdrojová forma práce ve formátu \LaTeX{}}.
		.1 text\DTcomment{text práce}.
		.2 thesis.pdf\DTcomment{text práce ve formátu PDF}.
	}
 % include `medium.tex' from `text/' subdirectory

\end{document}
